\section{Introduksjon}
I dette semesterprosjektet har vi laget en automatisk refill-arm, en robotarm som fyller glass med væske når en vektsensor registrerer at de er tomme. Prosjektet er gjennomført i emnet AUT-2606 ved Institutt for Automasjon og Prosessteknologi, og målet har vært å vise og videreutvikle vår kunnskap innen mekanismer, PCB-design og programmering.

Rapporten gir først en gjennomgang av det teoretiske grunnlaget for komponentene og metodene som er brukt, blant annet regulering og invers kinematikk. Deretter beskrives metodikken bak design av kretsen, lodding og programvaren, før vi presenterer og diskuterer resultatene fra det ferdige systemet.